% =====================================================================
%  Supplementary material: dataset generation
%  Three F3Loc-compatible floorplan-localization datasets with co-located
%  room impulse responses (Replica / Matterport3D / Structured3D).
%
%  Written to be cut down: every subsection is self-contained, so entire
%  blocks can be deleted without breaking the rest. Tables use booktabs.
%  Required packages: booktabs, amsmath, siunitx (optional), graphicx.
% =====================================================================

\section{Dataset generation}
\label{sec:supp-data}

We build three datasets that share one format, one coordinate convention and one
validation suite, so that a model trained on any of them can be evaluated on the
others without changes to the data loader. Each dataset pairs, at every camera
pose, (i) an RGB image, (ii) the ground-truth distance to the \emph{floor plan}
(walls only, furniture excluded), and (iii) a co-located active-acoustic
observation. Table~\ref{tab:supp-overview} summarises them.

\begin{table}[t]
\centering
\caption{Overview of the three datasets. A ``scene'' is one floor plan, i.e.\ one
occupancy map. For Matterport3D a building is split into one scene per storey.}
\label{tab:supp-overview}
\small
\begin{tabular}{lccc}
\toprule
 & Replica & Matterport3D & Structured3D \\
\midrule
Scenes (floor plans)            & 17 & 159 (83 buildings) & 700 \\
Split (train/val/test)          & 11/3/3 & 131/16/12 & 400/50/250 \\
Motion collections              & 2 & 2 & 1 \\
Views per sample $L{+}1$        & 4 & 4 & 1 \\
Frames per collection           & 20{,}400 & 65{,}780 & 15{,}872 \\
Samples per collection          & 5{,}100 & 16{,}445 & 15{,}872 \\
RIR poses per condition         & 5{,}100 & 16{,}445 & 15{,}872 \\
Acoustic conditions             & 2 & 2 & 1 \\
Image resolution                & $480{\times}640$ & $480{\times}640$ & $360{\times}640$ \\
Horizontal FOV                  & $106.3^\circ$ & $106.3^\circ$ & $80.0^\circ$ \\
$F_W = f/W$                     & $3/8$ & $3/8$ & $0.596$ \\
Total size                      & 46\,GB & 173\,GB & 30\,GB \\
\bottomrule
\end{tabular}
\end{table}

% ---------------------------------------------------------------------
\subsection{Directory layout}
\label{sec:supp-layout}

All three datasets use the same tree, so a loader parameterised by the dataset
root and the collection name reads any of them:

\begin{verbatim}
<dataset>/
  <collection>/ split.yaml
                <scene>/ rgb/                   RGB frames
                         poses.txt              x y yaw per frame
                         depth40.txt            40 ray-cast depths per frame
                         depth160.txt           160 ray-cast depths per frame
                         depth_radial_scan/     per-pixel radial depth, real geometry
                         depth_radial_floorplan/ per-pixel radial depth, floor plan
                         map.png                copy of the scene's occupancy map
                         chunks.json            per-frame metadata (see text)
  maps/<scene>/           map.png, scene_meta.json
  floorplan_proxy/<scene>/ floorplan.glb        extruded map, acoustic geometry
  desdf/<scene>/desdf.npy                       test scenes only
  rir/<collection>/<condition>/<scene>/pose_<i>/
        rir.npy, rir_metadata.json,
        rir_binaural.npy, rir_binaural_rel{090,180,270}.npy,
        rir_binaural_metadata.json
  dataset_meta.json                             machine-readable conventions
\end{verbatim}

Row order is positional: the $i$-th line of \texttt{poses.txt}, the $i$-th row of
\texttt{depth40.txt} and \texttt{depth160.txt}, the $i$-th file in ascending
filename order in \texttt{rgb/} and in both depth-map directories, and
\texttt{rir/.../pose\_\{i:05d\}} all describe the same frame.

% ---------------------------------------------------------------------
\subsection{Coordinate conventions}
\label{sec:supp-coords}

A pose is a \emph{global} $\mathrm{SE}(2)$ pose $(x, y, \theta)$ in the floor-plan
frame; poses are never relative to a previous frame. The origin is the centre of
\texttt{map.png}, $x$ and $y$ are metres, and $\theta$ is the heading measured
counter-clockwise from the $+x$ axis in radians, wrapped to $[-\pi, \pi]$.
Cameras are gravity aligned (zero roll and pitch), so a single scalar suffices
and no quaternion or rotation matrix is stored. World coordinates map to map
pixels as
\begin{equation}
  \mathrm{col} = \frac{x}{r} + \frac{W}{2}, \qquad
  \mathrm{row} = \frac{y}{r} + \frac{H}{2}, \qquad r = 0.01\,\mathrm{m/px},
\end{equation}
and the occupancy grid is indexed as \texttt{occ[row, col]}, which is the
convention obtained by displaying the map with \texttt{origin="lower"}.

The simulator (Habitat) is $y$-up while the floor plan is $z$-up. We use
\begin{equation}
  x = x_{h} - c_x, \qquad y = -z_{h} - c_y, \qquad \theta_{h} = \theta - \tfrac{\pi}{2},
\end{equation}
where $(c_x, c_y)$ is the map centre in simulator coordinates, stored per scene in
\texttt{scene\_meta.json}, and $\theta_h$ is the rotation about the simulator's
$+y$ axis. Sign errors here produce data that trains to a low loss and localises
at chance, so we verify the mapping numerically for every scene
(Sec.~\ref{sec:supp-validation}).

% ---------------------------------------------------------------------
\subsection{Floor-plan maps}
\label{sec:supp-maps}

\texttt{map.png} is an 8-bit PNG with three identical channels at
$0.01\,\mathrm{m/px}$, in which free space is \emph{exactly} 255 and every other
value is an obstacle. Walls \emph{and} the building exterior are obstacles. The
polarity is strict because the ray caster stops at the first pixel with
$255 - \mathrm{occ} > 0$, so an anti-aliased edge at 254 would be a wall; maps are
rasterised with nearest-neighbour and never passed through a lossy codec. One map
describes exactly one storey.

Crucially, the map is \emph{wall only}: furniture is excluded, because the
ground-truth depth is the distance to the floor plan and the network is trained to
see through clutter to structure. The three datasets reach that from different
sources.

\paragraph{Replica and Matterport3D (from the scan mesh).}
A wall is what is solid at \emph{every} height. Slicing the mesh at head height
would include sofas and tables; slicing just below the ceiling would include the
lintel above every door, which, once extruded, seals the doorway and destroys
room-to-room coupling. We therefore vote occupancy across $\approx 12$ height bands
spanning the storey, keeping only near-vertical faces ($|n_z| < 0.35$), and
threshold the vote at $0.6$: true walls score near $1.0$, whereas furniture (low
bands only) and lintels (high bands only) each reach only $\approx 0.3$. Scans have
holes (windows, unscanned corners), so the room mask alone does not enclose
anything; we intersect the navigation mesh with the scan footprint, fill it, and
seal its boundary with a one-cell ring of wall. Free space is the region reachable
by flood fill from the simulator's navigation graph nodes that lie inside that
envelope. The mask is built at $0.05\,\mathrm{m}$ and upsampled by 5 with
nearest-neighbour.

\paragraph{Matterport3D storeys.}
A Matterport3D scene is an entire building, which violates the one-map-per-floor
requirement. We cluster navigation-graph nodes by height (a gap larger than
$1\,\mathrm{m}$ starts a new storey) and emit one scene per storey, named
\texttt{<building>\_f<k>} with $k=0$ the lowest. All storeys of a building share a
single map frame; the vertical band of storey $k$ is
$[\,y_{\min}-0.1,\ \min(y_{\min}+2.7,\ y^{(k+1)}_{\min}-0.1)\,]$. Storeys with
fewer than four graph nodes, less than $8\,\mathrm{m}^2$ of sealed interior, or a
ceiling lower than the camera are discarded; the last criterion removes six
split-level bands that are too thin to hold a $1.25\,\mathrm{m}$ camera.

\paragraph{Structured3D (from structural annotations).}
Structured3D is synthetic and ships no mesh, but it provides the room polygons
directly. We rasterise the floor polygons on a $0.05\,\mathrm{m}$ grid with
4-connected wall lines, then re-open doors that connect two annotated rooms while
keeping exterior doors and windows solid: the map treats the exterior as obstacle
and the acoustic proxy has no exterior walls, so opening an entrance door would let
rays escape in the proxy but not on the map.

% ---------------------------------------------------------------------
\subsection{Camera and pose sampling}
\label{sec:supp-poses}

\paragraph{Camera.}
Replica and Matterport3D use the released F3Loc camera,
$K = [[240,0,320],[0,240,240],[0,0,1]]$ at $480{\times}640$, i.e.\ a horizontal
field of view of $2\arctan(320/240) \approx 106.3^\circ$ and $F_W = f/W = 3/8$. The
camera sits $1.25\,\mathrm{m}$ above the navigation mesh, co-located with the
acoustic source. Structured3D images are the renders shipped with that dataset,
downsized to $360{\times}640$; their field of view is $80^\circ$, giving
$F_W = 0.596$. \textbf{$F_W$ is not inferred from the images}: it must be set per
dataset in the network configuration, otherwise the ground-truth ray geometry and
the predicted rays disagree silently.

\paragraph{Motion regimes.}
The complementary network exists to arbitrate between well-conditioned and
degenerate baselines, so a dataset with a single motion regime cannot exercise it.
Replica and Matterport3D therefore ship two collections built from the same maps:

\begin{itemize}
  \item \textbf{Forward} (\texttt{*\_f}): each 4-frame sample translates
        $0.15$--$0.40\,\mathrm{m}$ per step along the heading with a yaw drift of
        $\mathcal{N}(0, 4^\circ)$ clipped to $\pm 10^\circ$.
  \item \textbf{General} (\texttt{*\_g}): 40\% forward ($\pm 20^\circ$ yaw,
        $0.10$--$0.40\,\mathrm{m}$), 35\% in-place rotation (monotonic
        $10$--$35^\circ$ per step, translation $\mathcal{N}(0, 2\,\mathrm{cm})$),
        25\% mixed ($0$--$0.30\,\mathrm{m}$ in any direction within $\pm 60^\circ$
        of the heading, yaw $\pm 35^\circ$). The in-place rotations are the
        degenerate case for the multi-view network.
\end{itemize}

Poses are drawn on the Habitat navigation mesh, so they are reachable and not
inside geometry, and are additionally required to lie on a free map pixel, to keep
at least $0.25\,\mathrm{m}$ of clearance from the nearest wall pixel, and to be
connected to the previous frame by a straight segment that does not cross a wall.
The simulator's sampler uses an internal random generator, so we re-seed it per
scene; without this every visit to a scene returns the same points. Replica uses
300 samples per scene; Matterport3D scales with storey area (0.6 per
$\mathrm{m}^2$, clipped to $[40, 160]$) because storey sizes span two orders of
magnitude. Structured3D poses are not sampled at all: we use the camera positions
shipped with the dataset, which is why $L = 0$ there. Structured3D cameras have
non-zero roll and pitch, so its images and depth maps are gravity-aligned with the
same homography used by F3Loc, and the original roll, pitch and camera height are
retained in \texttt{chunks.json}.

% ---------------------------------------------------------------------
\subsection{Ground-truth depth}
\label{sec:supp-depth}

\paragraph{Ray-cast structural depth.}
For $N \in \{40, 160\}$ we cast $N$ rays across the field of view,
\begin{equation}
  \alpha_i = \mathrm{flip}\!\left(\arctan \frac{i - \bar{i}}{N F_W}\right), \quad
  i = 0,\dots,N-1,
\end{equation}
and store the \emph{camera-forward} depth
$d_i = \rho(\theta + \alpha_i)\cos\alpha_i$, where $\rho$ is the range returned by
the ray caster on \texttt{map.png}. Storing range instead of forward depth is a
silent ${\approx}10\%$ error at the edges of the field of view. Rays that leave the
map saturate at $20\,\mathrm{m}$ and are kept finite; saturation is below $0.001\%$
of rays in all three datasets because the maps are sealed. \texttt{depth40}
supervises the monocular network and \texttt{depth160} the multi-view and
complementary networks.

\paragraph{Ray caster.}
We do not use the reference implementation's digital differential analyser. On
steps that move left or down it tests the pixel the ray is \emph{leaving} rather
than the one it enters, so a ray that clips the corner of a wall block passes
through it; along the staircase edge of a rotated wall this repeats at every block
and the ray escapes to the map border. We measured $0.27$--$0.43\%$ of rays
returning the saturation distance through a wall only $3\,\mathrm{m}$ away, and
DESDF inherits the same leaks. We therefore use an exact Amanatides--Woo traversal
that tests every pixel the ray passes through. On rays where the reference
implementation does not leak, the two agree to within one map pixel, and ours is
never longer.

\paragraph{Per-pixel radial depth.}
In addition to the 1-D structural depth required by the localisation networks, we
release dense depth at the image resolution, as 16-bit PNGs in millimetres with
$0$ marking no hit, in two variants: \texttt{depth\_radial\_scan} from the real
geometry (the scan mesh for Replica and Matterport3D, the shipped furnished render
for Structured3D) and \texttt{depth\_radial\_floorplan} from the extruded floor
plan. Values are \emph{radial}, i.e.\ the Euclidean distance along each pixel's
ray,
\begin{equation}
  d^{\mathrm{rad}}_{uv} = d^{z}_{uv}\sqrt{1 + \left(\tfrac{u + 0.5 - c_x}{f_x}\right)^2 + \left(\tfrac{v + 0.5 - c_y}{f_y}\right)^2},
\end{equation}
so that the difference between the two variants is exactly the furniture and the
scan geometry that the floor plan omits. Radial range is invariant to a rotation
about the camera centre, which is why the Structured3D depth is converted before
the gravity-alignment warp rather than after.

\paragraph{DESDF.}
For test scenes we precompute the directional ESDF used at evaluation time: a
dictionary \texttt{\{l, t, desdf\}} where \texttt{desdf} is
$(H, W, 36)$ \texttt{float32} in metres at $0.1\,\mathrm{m}$ per cell with 36
equiangular counter-clockwise orientations ($10^\circ$ per bin), cast to
$10\,\mathrm{m}$, and \texttt{l}, \texttt{t} are the crop offsets into
\texttt{map.png} ($x_{\mathrm{map}} = 10\,x_{\mathrm{desdf}} + l$). The crop is the
free-space bounding box, which keeps the array small.

% ---------------------------------------------------------------------
\subsection{Acoustic observation}
\label{sec:supp-acoustics}

At every evaluated pose (the reference frame of each sample) we render room impulse
responses with a compact array co-located with the camera: the agent emits and
listens at the same point, which makes the observation an active echolocation
measurement rather than a passive recording.

\paragraph{Array.}
Six omnidirectional microphones on a ring of radius $0.05\,\mathrm{m}$ at angles
$\{\pi, 4\pi/3, 5\pi/3, 2\pi, \pi/3, 2\pi/3\} + \theta$, with the simulator-frame
offsets $[\,r\cos a,\ 0,\ -r\sin a\,]$ (the negative sine is the $y = -z$ mapping of
Sec.~\ref{sec:supp-coords}). Channel 3 points along the heading. The renderer has
no multi-microphone array layout, so each channel is a separate mono render with
the receiver moved to that microphone and the source held at the ring centre. In
addition we render a two-channel binaural response with the simulator's built-in
head-related transfer function at four relative headings $0^\circ$, $90^\circ$,
$180^\circ$, $270^\circ$; the mono ring has no directivity, so rotating it would
merely resample the same circle and only the $0^\circ$ ring is stored.

\paragraph{Simulation parameters.}
Sample rate $48\,\mathrm{kHz}$; direct and indirect paths enabled; 20{,}000 indirect
rays with up to 50 reflections; edge diffraction enabled up to order 10;
transmission and per-surface materials disabled, so every surface takes the
renderer's default (absorption $0.10$, scattering $0.5$, frequency independent).
Impulse responses are $0.3$--$0.6\,\mathrm{s}$ long; the usable window is 6{,}144
samples starting 96 samples after the direct peak, which is the same
$2\,\mathrm{ms}$ guard and $128\,\mathrm{ms}$ window as the reference
specification expressed at $8\,\mathrm{kHz}$.

Three of these values depart from that specification and the reason is worth
stating. First, diffraction was specified off, but a dataset whose purpose is
non-line-of-sight inference cannot discard the mechanism that carries sound around
corners; enabling it costs nothing in the renderer. Second, the specified 4{,}096
rays leave the response dominated by sampling noise: rendering the same pose twice
with identical settings gives a waveform correlation of only $0.74$, rising to
$0.83$ at 16{,}384 rays and $0.91$ at 65{,}536; we use 20{,}000, which matches our
earlier renders. Third, the specified reflection limit of 6 truncates the response
at ${\approx}36\,\mathrm{ms}$, discarding exactly the higher-order paths that reach
an occluded receiver; 50 reflections reproduce the energy distribution of the
renderer's default limit of 200 at one third of the cost.

\paragraph{Conditions.}
Each pose is rendered against two geometries where both exist:
\texttt{raw\_scan\_open}, the scan mesh as captured, with its furniture, holes and
open boundaries; and \texttt{floorplan\_closed}, a watertight proxy obtained by
extruding the \emph{same} wall mask that produced \texttt{map.png} from floor to
ceiling, plus floor and ceiling slabs. The second condition exists because an
acoustic likelihood matched against a floor plan is only meaningful if the
acoustics and the floor plan describe the same room; the pair tells us whether an
audio-visual gain is real geometry or a mesh artefact. Structured3D has no mesh, so
only \texttt{floorplan\_closed} is available there; its \emph{visual} counterpart
of the scan condition is still present, since the shipped renders include
furniture.

% ---------------------------------------------------------------------
\subsection{Validation}
\label{sec:supp-validation}

Every scene passes an automated suite before release. Beyond bookkeeping checks
(row counts equal across modalities, sample sizes a multiple of $L{+}1$, filename
order equal to row order, disjoint splits, binary maps with free space exactly 255,
finite positive depths bounded by the saturation distance, DESDF present for every
test scene with the correct shape and dtype, and acoustic pose indices resolving to
the correct row), two geometric checks carry the load.

\paragraph{Axis check.}
The floor-plan proxy is the extruded map, so rendering it with the real pinhole
camera at a pose must reproduce, along the horizon row, exactly what the ray caster
returns on the map at the same pose and column angles. Any sign error, axis swap,
wrong $F_W$, or forward-versus-radial confusion appears here as tens of
centimetres. Across all three datasets the median absolute discrepancy is
$0.0\,\mathrm{cm}$ and the mean is below $0.1\,\mathrm{cm}$.

\paragraph{Independent re-derivation and full-frame survey.}
\texttt{depth160} is compared against a fine-step sampling march that shares no
code with the ray caster; the 99th percentile of the absolute difference is
$0.05\,\mathrm{cm}$, and the march is never shorter than the stored value. Finally,
\emph{every} frame's floor-plan depth map is scanned for no-hit pixels: because the
proxy is watertight and the map is sealed, there must be none. This check is what
surfaced two defects that no shape, range or NaN test would have caught: maps whose
sealing ring fell outside the grid, and cameras closer to a wall than the
renderer's default near plane.

\paragraph{Reported, not enforced.}
Scanned environments contain space that was never captured: Replica's furnished
apartments have unscanned volumes behind some doorways and no ceiling in several
scenes, and Matterport3D has similar gaps. A camera looking into such a void
returns a black image while the sealed map still reports a wall. We do not discard
those poses; instead every frame records the fraction of no-hit pixels in the lower
three quarters of the image in \texttt{chunks.json}, so a consumer can threshold it.
In Matterport3D the median over scenes is $3.8\%$ of frames above a $15\%$
threshold, with five of 318 scene--collection pairs above $90\%$.

% ---------------------------------------------------------------------
\subsection{Statistics}
\label{sec:supp-stats}

\begin{table}[t]
\centering
\caption{Per-scene statistics. Map sizes are in pixels at $0.01\,\mathrm{m/px}$;
free area is the sealed interior of one floor plan.}
\label{tab:supp-stats}
\small
\begin{tabular}{lccc}
\toprule
 & Replica & Matterport3D & Structured3D \\
\midrule
Free area (m$^2$), min / median / max & 9.4 / 36.5 / 75.5 & 19.1 / 164.1 / 2014.2 & 6.0 / 74.1 / 783.0 \\
Total mapped area (m$^2$)             & 768 & 39{,}432 & 53{,}868 \\
Map width (px), min / median / max    & 650 / 995 / 1505 & 700 / 2590 / 8365 & 450 / 1220 / 4050 \\
Frames per scene, min / median / max  & 1200 / 1200 / 1200 & 160 / 392 / 640 & 2 / 24 / 53 \\
Saturated depth rays                  & $0\%$ & $<0.01\%$ & $0\%$ \\
\bottomrule
\end{tabular}
\end{table}

\begin{table}[t]
\centering
\caption{Acoustic release. Impulse-response length varies per pose because the
renderer truncates once the energy has decayed.}
\label{tab:supp-acoustic}
\small
\begin{tabular}{lccc}
\toprule
 & Replica & Matterport3D & Structured3D \\
\midrule
Conditions                       & scan, floor plan & scan, floor plan & floor plan \\
Poses per condition              & 5{,}100 $\times$ 2 & 16{,}445 $\times$ 2 & 15{,}872 \\
Ring channels $\times$ headings  & $6 \times 1$ & $6 \times 1$ & $6 \times 1$ \\
Binaural channels $\times$ headings & $2 \times 4$ & $2 \times 4$ & $2 \times 4$ \\
Sample rate                      & \multicolumn{3}{c}{$48\,\mathrm{kHz}$} \\
Indirect rays / max reflections  & \multicolumn{3}{c}{20{,}000 / 50} \\
Diffraction order                & \multicolumn{3}{c}{10} \\
Guard / usable window (samples)  & \multicolumn{3}{c}{96 / 6{,}144} \\
\bottomrule
\end{tabular}
\end{table}

% ---------------------------------------------------------------------
\subsection{Limitations}
\label{sec:supp-limitations}

\begin{itemize}
  \item \textbf{No acoustic scan condition for Structured3D.} The dataset
        distributes renders and structural annotations but no mesh, so the
        furnished geometry does not exist in a form a ray tracer can use. Only the
        floor-plan condition is available there.
  \item \textbf{Single view for Structured3D.} We use the camera positions as
        shipped rather than sampling trajectories, so that subset trains the
        monocular branch only.
  \item \textbf{Unscanned space.} Replica and Matterport3D frames looking into
        unscanned volume are retained and flagged rather than removed; see
        Sec.~\ref{sec:supp-validation}.
  \item \textbf{Scene count.} Generalisation in this task is across floor plans,
        and Replica's 17 scenes are well below the 119 of the released reference
        data; Matterport3D (159 storeys) and Structured3D (700) are provided
        partly to compensate.
  \item \textbf{Materials.} All surfaces take the renderer's default acoustic
        material. Replica ships no material definitions, and using per-surface
        materials for one dataset only would make the three incomparable.
  \item \textbf{Thin storeys.} Six Matterport3D bands produced by the storey
        clustering were too thin to hold the camera and were discarded; a
        height-gap threshold tuned per building would recover some of them as
        split levels of the storey below.
\end{itemize}

% ---------------------------------------------------------------------
\subsection{Reproducibility}
\label{sec:supp-repro}

Each dataset ships \texttt{dataset\_meta.json}, a machine-readable record of every
convention above: the coordinate mapping, camera intrinsics, map resolution and
polarity, depth definition, DESDF layout, motion parameters, the full acoustic
configuration, the simulator version and the commit of the reference
implementation, the counts reported in Table~\ref{tab:supp-overview}, and an
explicit list of the points where we depart from the specification with the reason
for each. Per-scene records in \texttt{maps/<scene>/scene\_meta.json} carry the map
size, the origin offsets, the floor and ceiling heights, and the fraction of the
navigation mesh that falls on free map pixels. The generation and validation code
is released alongside the data; the entire pipeline is resumable, and rerunning a
stage regenerates only what is missing.
